\documentclass[12pt,a4paper]{article}

\usepackage[utf8]{inputenc}
\usepackage[T1]{fontenc}
\usepackage{lmodern}
\usepackage{textcomp}
\usepackage[french]{babel}
\frenchsetup{StandardLists=true,SmallCapsFigTabCaptions=false}
\usepackage[a4paper,margin=2.5cm,headheight=16pt]{geometry}
\usepackage{amsmath,amssymb}
\usepackage{xcolor}
\usepackage{graphicx}
\usepackage{booktabs}
\usepackage{multirow}
\usepackage{float}
\usepackage{needspace}
\usepackage{array}
\usepackage{enumitem}
\usepackage{fancyhdr}
\usepackage{titlesec}
\usepackage{caption}
\usepackage{siunitx}
\usepackage{tikz}
\usepackage[european,siunitx]{circuitikz}
\usetikzlibrary{arrows.meta,positioning,calc,babel}

% ---------- Couleurs ----------
\definecolor{bleu}{RGB}{0,67,130}
\definecolor{rouge}{RGB}{181,54,28}
\definecolor{fond}{RGB}{233,238,245}
\definecolor{gris}{RGB}{138,138,138}
\definecolor{noir}{RGB}{26,26,26}

\usepackage[colorlinks=true,linkcolor=bleu,citecolor=bleu,urlcolor=bleu]{hyperref}

% ---------- Unités ----------
\sisetup{output-decimal-marker={,},per-mode=power,detect-all}

% ---------- Titres ----------
\titleformat{\section}{\Large\bfseries\color{bleu}}{\thesection}{1em}{}
\titleformat{\subsection}{\large\bfseries\color{bleu}}{\thesubsection}{1em}{}
\titlespacing*{\section}{0pt}{3.0ex plus 1ex minus .2ex}{2.0ex plus .2ex}
\titlespacing*{\subsection}{0pt}{2.5ex plus 1ex minus .2ex}{1.5ex plus .2ex}

% ---------- Légendes ----------
\DeclareCaptionLabelSeparator{deuxpoints}{\NoAutoSpacing{:}\ }
\captionsetup{labelfont={bf,color=bleu},labelsep=deuxpoints,font=normalsize}
\captionsetup[table]{position=below}

% ---------- En-têtes ----------
\pagestyle{fancy}
\fancyhf{}
\fancyhead[L]{\itshape Compte-rendu – TP VNA}
\fancyhead[R]{\itshape ENSIL-ENSCI – Électronique \& Télécommunications}
\fancyfoot[C]{\thepage}
\renewcommand{\headrulewidth}{0.4pt}

\setlist[itemize]{label=\textbullet}
\setlength{\parindent}{0pt}
\setlength{\parskip}{0.6em}

\newcommand{\arel}{\textcolor{gris}{…}}

\begin{document}

% ================= Page de titre =================
\begin{center}
  ENSIL-ENSCI – Université de Limoges\\
  Spécialité Électronique \& Télécommunications – Travaux Pratiques Hyperfréquences\\[4pt]
  {\Large\bfseries\color{bleu} Compte-rendu – TP VNA}\\[6pt]
  {\bfseries Analyse vectorielle de réseaux – Application à un TEC}\\[6pt]
  Mouradi Souhail \ – Otmane Elbacha \ – Jawher TABBABI \quad – \quad Date : 16 octobre 2026\\[2pt]
  {\small\itshape La mise en page de ce document a été réalisée à l’aide de l’assistant IA Claude (Anthropic).}
\end{center}

\vspace{0.8em}
{\color{bleu}\rule{\linewidth}{1.2pt}}
\vspace{1.0em}

% ================= 1 =================
\section{Objectif}

L’objectif de ce TP est d’utiliser l’analyseur vectoriel de réseaux HP~8753C pour mesurer les
paramètres $S$ d’un transistor à effet de champ (TEC FUJITSU FLK012) dans la bande
\SIrange{2,95}{3,05}{\giga\hertz}. Il faut d’abord étalonner l’analyseur, puis polariser
correctement le TEC, et enfin comparer les mesures aux données du constructeur.

% ================= Matériel =================
\section{Matériel utilisé}

Le banc de mesure est présenté sur la Figure~\ref{fig:banc}. Chaque élément a un rôle précis :

\begin{itemize}
  \item \textbf{Analyseur vectoriel de réseaux HP~8753C.} C’est l’appareil de mesure principal.
        Il contient la source hyperfréquence (wobulateur) qui balaye la bande choisie, les
        récepteurs et l’écran. Il mesure le rapport entre l’onde réfléchie ou transmise et
        l’onde incidente, en module et en phase. Il affiche directement les paramètres $S_{ij}$
        (en dB, en degrés ou sur l’abaque de Smith). Il garde aussi en mémoire les termes
        d’erreur de l’étalonnage pour corriger les mesures.

  \item \textbf{Tiroir de test de paramètres $S$ HP~85047A.} Il est associé à l’analyseur et
        porte les deux ports de mesure. Il contient le commutateur qui envoie la source sur le
        port~1 ou sur le port~2, les coupleurs bidirectionnels qui prélèvent les ondes $a_i$
        et $b_i$, et la charge adaptée du port non utilisé. Grâce à lui, les quatre paramètres
        $S$ sont mesurés sans retourner le composant. Il possède aussi des entrées de
        polarisation qui amènent le continu jusqu’aux ports.

  \item \textbf{Câbles de mesure et kit d’étalonnage APC~\SI{7}{\milli\metre}.} Les câbles
        relient les ports au boîtier de test. Le kit contient les standards connus utilisés pour
        l’étalonnage SOLT : un court-circuit (Short), un circuit ouvert (Open) et une charge
        \SI{50}{\ohm} (Load). APC~7 est un connecteur de précision \SI{7}{\milli\metre}, sans
        sexe, très répétable : il convient bien aux mesures d’étalonnage.

  \item \textbf{Deux alimentations continues.} L’une fournit la tension de grille
        $V_{GS0} < 0$, l’autre la tension de drain $V_{DS0} > 0$, avec une masse commune.
        Leur limitation en courant protège le transistor : le courant de drain ne doit pas
        dépasser \SI{20}{\milli\ampere}.

  \item \textbf{Tés de polarisation.} Ils superposent le continu et le signal RF sur le même
        accès. Une capacité de découplage laisse passer la RF et bloque le continu vers
        l’analyseur. Une self de choc laisse passer le continu et empêche la RF de partir vers
        l’alimentation. Dans ce TP, ces éléments sont internes aux tés que nous utilisons.

  \item \textbf{Boîtier de test pour transistor HP~85041A.} Il reçoit le TEC, pastille en bas,
        grille côté port~1 et drain côté port~2. Les deux sources sont reliées à la masse du
        boîtier. Il fait le passage entre les connecteurs coaxiaux et les pattes du transistor.
        Ses accès ajoutent des pertes et un déphasage : c’est pourquoi un étalonnage TRL est
        ensuite nécessaire.

  \item \textbf{Kit TRL-Bis.} Ce sont des standards qui se placent dans le boîtier de test :
        THRU (avec cale), LINE (sans cale) et SHORT. Ils permettent de ramener les plans de
        référence aux bornes du transistor seul.

  \item \textbf{Adaptateur d’impédance.} Il permet de modifier l’impédance vue par le TEC
        pour l’adapter à \SI{50}{\ohm}. Il sert dans la suite du TP, pour réaliser
        l’adaptation du transistor et obtenir un amplificateur bande étroite.

  \item \textbf{Transistor FUJITSU FLK012.} C’est le composant mesuré : un TEC hyperfréquence
        en boîtier à quatre pattes (une grille, un drain et deux sources).
\end{itemize}

\begin{figure}[H]
\centering
\begin{tikzpicture}[font=\scriptsize,
  bloc/.style={draw,fill=fond,minimum height=0.7cm,align=center},
  fl/.style={-{Stealth[length=4pt]}}]
  \node[bloc,minimum width=3.4cm] (alim) at (0,2.6) {Alimentations continues\\$V_{GS0}$ et $V_{DS0}$};
  \node[bloc,minimum width=7.5cm,minimum height=1.0cm] (vna) at (0,1.0) {Analyseur HP~8753C + tiroir de test HP~85047A\\(source, commutateur, coupleurs, récepteurs)};
  \draw[fl] ($(alim.south)+(-0.8,0)$) -- ($(vna.north)+(-0.8,0)$) node[midway,left] {$V_{GS0}$};
  \draw[fl] ($(alim.south)+(0.8,0)$) -- ($(vna.north)+(0.8,0)$) node[midway,right] {$V_{DS0}$};
  \node[draw,fill=fond,thick,minimum width=2.6cm,minimum height=1.1cm,align=center] (tec) at (0,-1.4) {Boîtier HP~85041A\\TEC FLK012};
  \draw[thick] ($(vna.south)+(-2.8,0)$) |- node[pos=0.25,left,align=right] {Port 1\\(grille)} (tec.west);
  \draw[thick] ($(vna.south)+(2.8,0)$) |- node[pos=0.25,right,align=left] {Port 2\\(drain)} (tec.east);
  \draw[rouge,densely dashed] ($(tec.west)+(-0.15,-0.8)$) -- ($(tec.west)+(-0.15,0.8)$);
  \draw[rouge,densely dashed] ($(tec.east)+(0.15,-0.8)$) -- ($(tec.east)+(0.15,0.8)$);
  \node[rouge] at (0,-2.35) {Plans de référence};
  \node[gris,font=\scriptsize\itshape] at (0,-2.75) {Câbles APC~7 entre les ports et le boîtier};
\end{tikzpicture}
\caption{Banc de mesure du TEC non adapté dans son boîtier de test.}
\label{fig:banc}
\end{figure}

% ================= 2 =================
\section{Partie théorique}

\textbf{Paramètres $S$.} Un quadripôle est décrit par les ondes incidentes $a_i$ et réfléchies
$b_i$ sur chaque accès. Elles sont liées par la matrice $S$ :
\begin{equation}
  \begin{pmatrix} b_1 \\ b_2 \end{pmatrix} =
  \begin{pmatrix} S_{11} & S_{12} \\ S_{21} & S_{22} \end{pmatrix}
  \begin{pmatrix} a_1 \\ a_2 \end{pmatrix}
  \qquad\text{avec}\qquad
  S_{ij} = \left.\frac{b_i}{a_j}\right|_{a_k = 0,\ k \neq j}
  \label{eq:S}
\end{equation}

\textbf{Principe de l’analyseur.} Un wobulateur interne excite le dispositif sous test (DST) sur
l’accès~1 ou l’accès~2, selon la position d’un commutateur. L’accès non utilisé est fermé sur une
charge adaptée \SI{50}{\ohm}, ce qui impose $a_2 = 0$ ou $a_1 = 0$. Deux coupleurs
bidirectionnels prélèvent des ondes $a_{iM}$ et $b_{jM}$ proportionnelles à $a_i$ et $b_j$
(Figure~\ref{fig:vna}). Après changement de fréquence (mélangeurs), on obtient deux signaux BF
$x(t)$ et $y(t)$ : le rapport de leurs amplitudes donne $|b_j/a_i|$ et leur déphasage donne
$\arg(b_j/a_i)$.

\begin{figure}[H]
\centering
\begin{tikzpicture}[font=\scriptsize,
  bloc/.style={draw,fill=fond,minimum height=0.6cm,align=center},
  fl/.style={-{Stealth[length=4pt]}},
  onde/.style={-{Stealth[length=4pt]},bleu,thick}]
  % Générateur et commutateur
  \draw (0,3.3) circle (0.28);
  \draw (-0.15,3.3) sin (-0.07,3.4) cos (0,3.3) sin (0.07,3.2) cos (0.15,3.3);
  \node[right=6pt] at (0.28,3.3) {Wobulateur};
  \draw (0,3.02) -- (0,2.6);
  \draw[thick] (0,2.6) -- (-0.5,2.25);
  \fill (0,2.6) circle (1.2pt);
  \node[right] at (0.1,2.45) {Commutateur};
  \draw (-0.5,2.2) -- (-3.6,2.2) -- (-3.6,0.25);
  \draw[dashed] (0.5,2.2) -- (3.6,2.2) -- (3.6,0.25);
  % charge adaptée
  \draw (0.5,2.2) -- (0.5,1.9);
  \draw[fill=white] (0.38,1.9) rectangle (0.62,1.45);
  \draw (0.5,1.45) -- (0.5,1.3);
  \draw (0.3,1.3) -- (0.7,1.3);
  \draw (0.37,1.22) -- (0.63,1.22);
  \draw (0.44,1.14) -- (0.56,1.14);
  \node[right,align=left] at (0.65,1.62) {Charge\\[-2pt]\SI{50}{\ohm}};
  % Coupleurs et DST
  \node[bloc,minimum width=1.6cm] (c1) at (-3.6,0) {Coupleur};
  \node[bloc,minimum width=1.6cm] (c2) at (3.6,0) {Coupleur};
  \node[above left=1pt and -2pt of c1.north west] {Port 1};
  \node[above right=1pt and -2pt of c2.north east] {Port 2};
  \node[draw,fill=fond,minimum width=1.2cm,minimum height=1.3cm,align=center,thick] (dst) at (0,0) {DST\\(TEC)};
  \draw (c1.east) -- (dst.west);
  \draw (dst.east) -- (c2.west);
  \draw[onde] (-2.0,0.18) -- (-1.0,0.18) node[midway,above] {$a_1$};
  \draw[onde] (-1.0,-0.18) -- (-2.0,-0.18) node[midway,below] {$b_1$};
  \draw[onde] (2.0,0.18) -- (1.0,0.18) node[midway,above] {$a_2$};
  \draw[onde] (1.0,-0.18) -- (2.0,-0.18) node[midway,below] {$b_2$};
  \draw[rouge,densely dashed] (-0.85,-1.0) -- (-0.85,1.0);
  \draw[rouge,densely dashed] (0.85,-1.0) -- (0.85,1.0);
  \node[rouge] at (0,-1.2) {Plans de référence};
  % Boîte de commutation
  \node[bloc,minimum width=9cm] (bc) at (0,-2.3) {Boîte de commutation};
  \draw[fl] ($(c1.south)+(-0.4,0)$) -- ($(c1.south)+(-0.4,-1.7)$) node[pos=0.45,left] {$a_{1M}$};
  \draw[fl] ($(c1.south)+(0.4,0)$) -- ($(c1.south)+(0.4,-1.7)$) node[pos=0.45,right] {$b_{1M}$};
  \draw[fl] ($(c2.south)+(-0.4,0)$) -- ($(c2.south)+(-0.4,-1.7)$) node[pos=0.45,left] {$a_{2M}$};
  \draw[fl] ($(c2.south)+(0.4,0)$) -- ($(c2.south)+(0.4,-1.7)$) node[pos=0.45,right] {$b_{2M}$};
  % Mélangeurs
  \node[draw,circle,fill=white,inner sep=0pt,minimum size=0.55cm] (m1) at (-2.2,-3.8) {$\times$};
  \node[draw,circle,fill=white,inner sep=0pt,minimum size=0.55cm] (m2) at (2.2,-3.8) {$\times$};
  \draw[fl] (-2.2,-2.6) -- (m1.north) node[midway,left] {$b_{jM}$};
  \draw[fl] (2.2,-2.6) -- (m2.north) node[midway,right] {$a_{iM}$};
  \draw (m1.east) -- (m2.west) node[midway,above] {OL};
  \node[left=2pt of m1] {Mélangeur};
  \node[right=2pt of m2] {Mélangeur};
  % Traitement BF
  \node[bloc,minimum width=6cm] (bf) at (0,-5.2) {Récepteur FI : rapport d’amplitude $X/Y$ et déphasage $\Delta\Phi$};
  \draw[fl] (m1.south) -- (-2.2,-4.9) node[midway,left] {$x(t)$};
  \draw[fl] (m2.south) -- (2.2,-4.9) node[midway,right] {$y(t)$};
  \node[bloc,minimum width=3cm] (af) at (0,-6.3) {Affichage $S_{ij}$};
  \draw[fl] (bf.south) -- (af.north);
  \node[bleu,font=\scriptsize\itshape] at (0,-7.0) {Flèches bleues : ondes de puissance dans les plans de référence};
\end{tikzpicture}
\caption{Synoptique simplifié de l’analyseur vectoriel de réseaux.}
\label{fig:vna}
\end{figure}

\textbf{Étalonnage.} Les coupleurs, les câbles et le commutateur ne sont pas parfaits (directivité
finie, désadaptation, fuites). Ces défauts sont décrits par un modèle d’erreur à 12 termes : 6 par
sens de mesure (directivité, adaptation de source, réponse en réflexion, réponse en transmission,
adaptation de charge, isolation). On les détermine en mesurant des standards connus :
\begin{itemize}
  \item \textbf{SOLT} (\textit{Short, Open, Load, Thru}) : court-circuit, circuit ouvert,
        charge \SI{50}{\ohm} et liaison directe. Les plans de référence sont au bout des câbles ;
  \item \textbf{TRL} (\textit{Thru, Reflect, Line}) : réalisé avec des standards dans la même
        technologie que le boîtier. Il ramène les plans de référence aux bornes du transistor seul.
\end{itemize}
L’étalonnage n’est valable que dans la bande de fréquence et pour la puissance choisies.

\textbf{TEC.} Le courant de drain $I_D$ est commandé par la tension grille-source $V_{GS}$
($-V_P < V_{GS} < 0$). Sur les caractéristiques $I_D(V_{DS})$, on distingue la zone ohmique, la
zone saturée (amplification), la zone d’avalanche et la zone de conduction de la diode
grille. Le point de repos doit être dans la zone saturée. Une incursion dans la zone
d’avalanche détruit le transistor.

\textbf{Stabilité et gain.} À partir des paramètres $S$, on calcule le facteur de stabilité :
\begin{equation}
  K = \frac{1 - |S_{11}|^2 - |S_{22}|^2 + |\Delta|^2}{2\,|S_{12}\,S_{21}|}
  \qquad\text{avec}\qquad
  \Delta = S_{11}S_{22} - S_{12}S_{21}
  \label{eq:K}
\end{equation}
Le TEC est inconditionnellement stable si $K > 1$ et $|\Delta| < 1$. Le gain maximum est alors
$G_{\max} = \frac{|S_{21}|}{|S_{12}|}\bigl(K - \sqrt{K^2 - 1}\bigr)$. Sinon, on utilise le gain
maximum stable $\mathrm{MSG} = |S_{21}|/|S_{12}|$.

% ================= 3 =================
\section{Étalonnage SOLT de l’analyseur}

L’analyseur HP~8753C est associé au tiroir de test HP~85047A. Les réglages sont donnés dans le
Tableau~\ref{tab:reglages}.

\begin{table}[H]
\centering
\begin{tabular}{ll}
\toprule
\textbf{Réglage} & \textbf{Valeur} \\
\midrule
Bande de fréquence          & \SIrange{2,95}{3,05}{\giga\hertz} \\
Puissance de la source      & \SI{-10}{dBm} \\
Kit d’étalonnage            & APC \SI{7}{\milli\metre} \\
Type d’étalonnage           & \textit{Full Two Ports} \\
\bottomrule
\end{tabular}
\caption{Réglages de l’analyseur avant l’étalonnage.}
\label{tab:reglages}
\end{table}

La puissance est réglée à \SI{-10}{dBm} \textbf{avant} l’étalonnage. Elle doit être faible pour
que le TEC reste en régime petit signal, donc linéaire. Si on la change après, l’étalonnage n’est
plus valable.

L’étalonnage a été fait dans l’ordre suivant :
\begin{itemize}
  \item \textbf{Réflexion} : Short, Open et Load sur le port~1, puis sur le port~2 ;
  \item \textbf{Transmission} : les deux câbles connectés ensemble (\textit{Thru}) ;
  \item \textbf{Isolation} : les deux ports fermés sur une charge \SI{50}{\ohm}.
\end{itemize}

\subsection{Vérification de l’étalonnage}

Pour valider l’étalonnage, on mesure des éléments dont on connaît la réponse
(Tableau~\ref{tab:verif}).

\begin{table}[H]
\centering
\begin{tabular}{llccc}
\toprule
\textbf{Montage} & \textbf{Paramètre} & \textbf{Module attendu} & \textbf{Phase attendue} & \textbf{Résultat} \\
\midrule
Câbles connectés  & $S_{21}$, $S_{12}$ & \SI{0}{dB}   & \SI{0}{\degree}   & Conforme \\
(\textit{Thru})   & $S_{11}$, $S_{22}$ & très faible  & –                 & Conforme \\
\midrule
Court-circuit     & $S_{11}$, $S_{22}$ & \SI{0}{dB}   & \SI{180}{\degree} & Conforme \\
\bottomrule
\end{tabular}
\caption{Vérification de l’étalonnage SOLT sur toute la bande.}
\label{tab:verif}
\end{table}

Avec les câbles connectés, l’onde passe entièrement d’un port à l’autre sans déphasage :
$|S_{21}| = |S_{12}| = \SI{0}{dB}$ et la phase est nulle. Il n’y a presque pas de réflexion.
Avec le court-circuit, toute l’onde est réfléchie avec un signe opposé ($\Gamma = -1$) :
$|S_{11}| = \SI{0}{dB}$ et la phase vaut \SI{180}{\degree}. Les traces obtenues sont plates sur
toute la bande et correspondent aux valeurs attendues : l’étalonnage est validé. On peut donc
passer à la mesure du TEC.

% ================= 4 =================
\needspace{10\baselineskip}
\section{Mesure du TEC FLK012}

\subsection{Montage et polarisation}

Le TEC est placé dans le boîtier de test HP~85041A, pastille en bas. La grille est reliée au
port~1, le drain au port~2, et les deux sources sont à la masse du boîtier. La polarisation
continue est amenée par deux tés de polarisation (Figure~\ref{fig:polar}). La capacité laisse
passer le signal RF et bloque le continu. La self bloque la RF et laisse passer le continu.

\begin{figure}[H]
\centering
\begin{circuitikz}[scale=0.85,transform shape,font=\small]
  % Té grille
  \draw[densely dashed,bleu,fill=fond] (0.9,-0.7) rectangle (3.3,2.2);
  \node[bleu] at (2.1,-1.0) {Té de polarisation};
  \draw (0,0) node[left,align=right]{Entrée RF\\(port 1)} to[short,o-] (1,0) to[C] (2.6,0) coordinate(g);
  \draw (g) to[L,*-] (2.6,2.0) -- (2.6,2.6) node[above]{$-V_P < V_{GS0} < 0$};
  \draw (g) -- (4.6,0);
  % TEC
  \node[njfet,anchor=G] (Q) at (4.6,0) {};
  \node[above] at ($(Q.G)+(-0.35,0)$) {G};
  \node[left] at ($(Q.D)+(0,-0.15)$) {D};
  \node[left] at ($(Q.S)+(0,0.15)$) {S};
  \draw (Q.S) node[ground]{};
  % Té drain
  \draw[densely dashed,bleu,fill=fond] (6.5,0.6) rectangle (8.9,3.5);
  \node[bleu] at (7.7,0.3) {Té de polarisation};
  \draw (Q.D) |- (7.2,1.3) coordinate(d);
  \draw (d) to[L,*-] (7.2,3.3) -- (7.2,3.9) node[above]{$V_{DS0} > 0$};
  \draw (d) to[C] (8.8,1.3) to[short,-o] (9.6,1.3) node[right,align=left]{Sortie RF\\(port 2)};
\end{circuitikz}
\caption{Polarisation du TEC à travers deux tés de polarisation.}
\label{fig:polar}
\end{figure}

Le point de repos choisi est en classe~A (Tableau~\ref{tab:polar}). La limitation en courant des
alimentations est réglée pour que le courant de drain ne dépasse pas \SI{20}{\milli\ampere}.

\begin{table}[H]
\centering
\begin{tabular}{cccc}
\toprule
$V_{DS0}$ & $I_{D0}$ & $V_{GS0}$ & $I_{D,\max}$ \\
\midrule
\SI{3}{\volt} & \SI{10}{\milli\ampere} & \SI{-0,4}{\volt} & \SI{20}{\milli\ampere} \\
\bottomrule
\end{tabular}
\caption{Point de polarisation en classe A.}
\label{tab:polar}
\end{table}

L’ordre des opérations est important pour ne pas détruire le transistor :
\begin{itemize}
  \item \textbf{polarisation} : vérifier les polarités ($V_{GS} < 0$, $V_{DS} > 0$, masse
        commune), régler les deux tensions à \SI{0}{\volt}, connecter la grille puis le drain,
        amener $V_{GS0}$ à \SI{-0,4}{\volt}, puis monter $V_{DS0}$ jusqu’à \SI{3}{\volt} ;
  \item \textbf{dépolarisation} : descendre $V_{DS0}$ puis $V_{GS0}$ à \SI{0}{\volt},
        déconnecter le drain et la grille, puis couper les alimentations.
\end{itemize}
Il ne faut jamais augmenter $V_{DS0}$ quand $V_{GS0} = \SI{0}{\volt}$ : l’avalanche est alors
brutale et irréversible.

\subsection{Mesure des paramètres $S$}

L’étalonnage SOLT donne les paramètres $S$ du TEC aux bornes du boîtier : ils incluent les
accès du boîtier. Pour obtenir les paramètres du transistor seul, un étalonnage TRL est fait avec
le kit TRL-Bis dans la même bande. Les standards utilisés sont : THRU (avec cale), LINE (sans
cale) et SHORT. Les relevés à \SI{3}{\giga\hertz} seront regroupés dans le
Tableau~\ref{tab:S}.

\begin{table}[H]
\centering
\begin{tabular}{c cc cc cc}
\toprule
 & \multicolumn{2}{c}{\textbf{Boîtier (SOLT)}} & \multicolumn{2}{c}{\textbf{Transistor (TRL)}} & \multicolumn{2}{c}{\textbf{Constructeur}} \\
\cmidrule(lr){2-3}\cmidrule(lr){4-5}\cmidrule(lr){6-7}
\textbf{Paramètre} & $|S|$ (dB) & $\varphi$ (°) & $|S|$ (dB) & $\varphi$ (°) & $|S|$ (dB) & $\varphi$ (°) \\
\midrule
$S_{11}$ & \arel & \arel & \arel & \arel & \arel & \arel \\
$S_{21}$ & \arel & \arel & \arel & \arel & \arel & \arel \\
$S_{12}$ & \arel & \arel & \arel & \arel & \arel & \arel \\
$S_{22}$ & \arel & \arel & \arel & \arel & \arel & \arel \\
\bottomrule
\end{tabular}
\caption{Paramètres $S$ du TEC à \SI{3}{\giga\hertz}, $V_{DS0} = \SI{3}{\volt}$, $I_{D0} = \SI{10}{\milli\ampere}$.}
\label{tab:S}
\end{table}

On s’attend à $|S_{21}| > \SI{0}{dB}$ (le TEC amplifie) et à $|S_{12}|$ très faible (le TEC est
presque unilatéral). Les écarts entre SOLT et TRL viennent des accès du boîtier, qui ajoutent
des pertes et un déphasage. Les valeurs TRL doivent être les plus proches de celles du
constructeur. Les mesures permettront ensuite de calculer $K$ et $|\Delta|$ avec
l’équation~(\ref{eq:K}).

% ================= 5 =================
\section{Conclusion}

L’analyseur a été étalonné en SOLT \textit{Full Two Ports} dans la bande
\SIrange{2,95}{3,05}{\giga\hertz}, avec une puissance de \SI{-10}{dBm}. La vérification avec les
câbles connectés ($\SI{0}{dB}$, \SI{0}{\degree}) et avec le court-circuit ($\SI{0}{dB}$,
\SI{180}{\degree}) donne les valeurs attendues : l’étalonnage est validé.

La suite du TP porte sur la polarisation du TEC FLK012 en classe~A ($V_{DS0} = \SI{3}{\volt}$,
$I_{D0} = \SI{10}{\milli\ampere}$), puis sur la mesure de ses paramètres $S$ et leur comparaison
avec la notice du constructeur.

\end{document}
